% ─────────────────────────────────────────────────────────────────────────────
\section{Glossary}
\label{sec:glossary}
% ─────────────────────────────────────────────────────────────────────────────

\begin{description}[style=standard,leftmargin=4.5cm,labelwidth=4.3cm,font=\ttfamily\bfseries\small]

  \item[Activation function]
    A non-linear function applied to the weighted sum of inputs at each neural network
    node.  MoonAI supports \emph{Sigmoid}, \emph{Tanh}, and \emph{ReLU}.

  \item[Adjusted fitness]
    A genome's raw fitness divided by its species size (explicit fitness sharing).
    Prevents any single large species from dominating reproduction.

  \item[AgentId]
    \texttt{using AgentId = std::uint32\_t}.  Unique identifier assigned to each agent
    at generation start; stable within one generation.

  \item[CSR]
    Compressed Sparse Row — a memory layout that packs variable-length per-agent network
    data into flat arrays with offset descriptors (\texttt{GpuNetDesc}), eliminating
    VRAM waste from padding.

  \item[CUDA]
    NVIDIA's parallel computing platform.  MoonAI uses CUDA to accelerate NN inference
    (\texttt{neural\_forward\_kernel}) and fitness evaluation
    (\texttt{fitness\_eval\_kernel}).

  \item[eval\_order]
    The topological sort of a network's nodes produced by Kahn's algorithm at
    \texttt{NeuralNetwork} construction time.  Stored as a \texttt{vector<uint32\_t>}
    of node ids in activation order.

  \item[Fast-forward mode]
    An optional rendering mode (\textbf{H} key) where the visualization skips drawing
    frames, allowing the simulation to run at maximum CPU/GPU throughput.  The HUD
    displays \texttt{[FF]} when active.

  \item[FPS]
    Frames Per Second — the rendering rate measured by \texttt{VisualizationManager}.
    Target set by \texttt{target\_fps} in config.

  \item[GpuBatch]
    RAII class that owns all CUDA device memory for one generation's population.
    Allocates on construction; frees via \texttt{cudaFree} in destructor.

  \item[GpuNetDesc]
    Eight-integer struct describing one agent's network within the CSR flat arrays
    (offsets: \texttt{node\_off}, \texttt{eval\_off}, \texttt{conn\_off},
    \texttt{out\_off}; counts: \texttt{num\_nodes}, \texttt{num\_eval},
    \texttt{num\_inputs}, \texttt{num\_outputs}).

  \item[Golden angle]
    $137.508^\circ = 360^\circ(1 - 1/\phi)$ where $\phi$ is the golden ratio.
    Used by \texttt{Renderer::species\_color()} to generate maximally distinct hues
    for consecutive species IDs.

  \item[HLD]
    High-Level Design document (December 2025) — conceptual architecture preceding this
    LLD.

  \item[IEEE 1016-2009]
    IEEE Standard for Software Design Descriptions.  This document follows the structure
    defined therein.

  \item[InnovationTracker]
    Per-generation map from connection pairs \texttt{(in\_node, out\_node)} to
    innovation numbers.  Guarantees that the same structural mutation receives the same
    historical marking within a generation, enabling correct NEAT crossover.

  \item[LLD]
    Low-Level Design — this document.

  \item[MT19937]
    Mersenne Twister 19937 — the pseudorandom number generator algorithm.  MoonAI uses
    the 64-bit variant \texttt{std::mt19937\_64} via the \texttt{Random} class.

  \item[NEAT]
    NeuroEvolution of Augmenting Topologies — the evolutionary algorithm that
    simultaneously evolves neural network weights and topology.  Reference: Stanley \&
    Miikkulainen (2002)~\cite{stanley2002}.

  \item[NN]
    Neural Network — a \texttt{NeuralNetwork} object built from a \texttt{Genome}.

  \item[RAII]
    Resource Acquisition Is Initialization — C++ design idiom ensuring deterministic
    resource lifetime via constructors and destructors.

  \item[RNG]
    Random Number Generator.  In MoonAI, one \texttt{Random} instance is shared per
    simulation for full reproducibility.

  \item[SensorInput]
    15-element observation vector produced by \texttt{Physics::build\_sensors()}.
    Encodes distances and angles to nearest predator, prey, and food; energy level;
    velocity components; density counts; and wall proximity.

  \item[SFML]
    Simple and Fast Multimedia Library — used for window creation, 2-D rendering, and
    event handling in the visualization package.

  \item[SpatialGrid]
    Hash-grid partitioning the world into cells of size \texttt{vision\_range/2}.
    Enables O(1) amortized neighbor queries for sensor building and attack detection.

  \item[Species stagnation]
    A species that has shown no improvement in best-ever fitness for
    \texttt{stagnation\_limit} consecutive generations is removed by
    \texttt{remove\_stagnant\_species()}, except the current top-fitness species.

  \item[TickCallback]
    \texttt{std::function<void(int, const SimulationManager\&)>} — optional hook
    installed on \texttt{EvolutionManager} to receive per-tick notifications
    (for logging, visualization updates, and benchmarking).

  \item[Topological sort]
    Total ordering of a DAG's nodes such that every edge $(u,v)$ has $u$ before $v$.
    Computed by Kahn's algorithm in \texttt{NeuralNetwork} constructor to determine
    activation evaluation order.

  \item[UML]
    Unified Modeling Language.  Version 2.5 class and package diagrams are used
    throughout this document.

  \item[Vec2]
    \texttt{struct Vec2 \{ float x, y; \}} — the 2-D coordinate type used everywhere in
    simulation physics.

  \item[VRAM]
    Video RAM — device memory on the GPU.  CSR layout minimizes VRAM usage relative to
    padded alternatives.

\end{description}
